\chapter{Regionalligen: Kopplung, Aufstieg, Malus}\label{kap:regionalliga}

In den Regionalligen hängt an einem Tabellenplatz oft nicht allein, ob man
auf- oder absteigt. Wie viele Mannschaften eine Staffel abgeben muss, hängt
davon ab, wie viele Drittligisten in sie zurückfallen. Wer aufsteigt, hängt
für zwei Staffeln an Aufstiegsspielen gegeneinander. Die Rechtsgrundlagen
für 2026/27 stehen in \datei{docs/abstieg_aufstieg_RL_2026_2027.md}, die
Modellannahmen in \datei{docs/modellannahmen-rl-kopplung.md}. Die
Entscheidungen dazu halten ADR~0006 und ADR~0008 fest.

\section{Die Abstiegskopplung an die 3.~Liga}\label{sec:kopplung}

\subsection*{Verteilung der Drittliga-Absteiger je Staffel}

Jeder Drittligist ist über die Spalte \code{Region} der Stammdatei einer
Staffel zugeordnet. Aus der Zusatzzählung des Monte-Carlo-Laufs
(Kapitel~\ref{kap:simulation}) folgt direkt die Verteilung der Zahl der
Absteiger je Staffel. Sie ist ausgezählt und nicht genähert
(\datei{RCode/rl_abstiegskopplung.R:116--181}):
\begin{equation}
  q_g(j) = P(\text{genau } j \text{ Drittliga-Absteiger in Staffel } g) = \frac{G_{g,j}}{N}.
\end{equation}
Ausgezählt und nicht etwa aus den einzelnen Abstiegswahrscheinlichkeiten
zusammengesetzt wird, weil die Abstiegsplätze negativ korreliert sind:
Steigt ein Verein aus dem Norden ab, ist ein Platz weniger für alle anderen
frei. Der Code prüft zwei Invarianten. Jede Zeile summiert sich zu $N$, und
die erwartete Gesamtzahl $\sum_g \sum_j j\,q_g(j)$ ist exakt $K = 4$.

\subsection*{Abstiegsplätze je Staffel}

Wie viele Plätze bei $j$ Drittliga-Absteigern zum Abstieg führen, regeln die
Spielordnungen der Regionalverbände (\datei{rl_abstiegskopplung.R:65--86}):

\begin{center}
\begin{tabular}{@{}llp{6.2cm}@{}}
\toprule
Staffel & Abstiegsplätze $n_g(j)$ & Grundlage \\
\midrule
Nord     & $3 + j$             & NFV-SpO §\,6 Abs.~3/4 \\
Nordost  & $\min(1 + j,\ 2)$   & NOFV A.\ Nr.~5, Deckel 2 als Modellannahme \\
West     & $4$                 & WDFV, entkoppelt \\
SüdWest  & $\min(3 + j,\ 5)$   & RLSW-SpO §\,47 Nr.~1/2 \\
Bayern   & $2$                 & BFV, dazu Relegation (s.\,u.) \\
\bottomrule
\end{tabular}
\end{center}

\subsection*{Platzgewichte und Abstiegswahrscheinlichkeit}

Für eine Staffel mit $n$ Mannschaften ist der $d$-letzte Platz genau dann ein
Abstiegsplatz, wenn $n_g(j) \ge d$. Sein Gewicht ist die Wahrscheinlichkeit
dafür (\datei{rl_abstiegskopplung.R:339--345}):
\begin{equation}
  w_g(n - d + 1) = \sum_{j \,:\, n_g(j) \ge d} q_g(j), \qquad d = 1, \dots, n.
\end{equation}
Die Abstiegswahrscheinlichkeit einer Mannschaft verbindet dieses Gewicht mit
ihrer eigenen Platzverteilung aus dem Hauptlauf der Staffel
(\datei{rl_abstiegskopplung.R:371--392}):
\begin{equation}\label{eq:rl-abstieg}
  \boxed{\;P(t \text{ steigt ab}) = \sum_{p=1}^{n} \hat P_{t,p}\; w_g(p)\;}
\end{equation}
Das Produkt ist exakt und keine Näherung. 3.~Liga und Regionalliga sind
getrennte Wettbewerbe, ihre Ergebnisse also unabhängig.

\paragraph{Nord: die zweite Kopplung.} Steigt der Nord-Meister auf, muss die
Staffel einen Absteiger weniger stellen. Mit $p_M$ als Summe der
Aufstiegswahrscheinlichkeiten der Staffel mischt der Code beide Fälle
(\datei{RCode/update_all_leagues_loop.R:775--780}):
\begin{equation}
  w^{\text{Nord}} = \clamp_{[0,1]}\Bigl((1 - p_M)\, w[n_g(j)] + p_M\, w[\max(n_g(j) - 1,\ 0)]\Bigr).
\end{equation}
Hier handelt es sich ausdrücklich um eine Näherung: Meisterschaft und
Abstieg stammen aus derselben Simulation der Staffel und sind nicht
unabhängig.

\paragraph{Bayern: Relegation.} Die beiden Letzten steigen direkt ab. Die
Plätze $n-3$ und $n-2$ spielen eine Relegation gegen die Bayernligen, die das
Modell nicht auflöst. Die Website weist diese Plätze als eigene Spalte
\enquote{Relegation} aus:
$P(\text{Relegation}) = \hat P_{t,n-3} + \hat P_{t,n-2}$.

\section{Wer aufsteigt}

Vier Regionalligisten steigen in die 3.~Liga auf. West und SüdWest stellen
ihren Aufsteiger dauerhaft direkt. Der dritte direkte Platz rotiert unter
Nord, Nordost und Bayern. Die beiden übrigen dieser drei ermitteln den
vierten Aufsteiger in Hin- und Rückspiel. Wer in welcher Saison an der Reihe
ist, steht als Daten in \code{AUFSTIEGSROTATION}
(\datei{RCode/rl_aufstieg.R:53}). Für 2026/27 steigt Nordost direkt auf,
\textbf{Nord und Bayern} spielen gegeneinander. Für eine unbekannte Saison
bricht der Code ab, statt still mit dem Vorjahreswert weiterzurechnen.

\subsection*{Direktaufsteiger}

\begin{equation}
  P(\text{Aufstieg}) = P(\text{Meister}) = \hat P_{t,1}.
\end{equation}

\subsection*{Aufstiegsspiele}

Für die beiden Playoff-Staffeln summiert der Code über alle möglichen
Meisterpaarungen. Weil die Staffeln unabhängig simuliert werden, ist auch
diese Doppelsumme exakt (\datei{rl_aufstieg.R:120--190}):
\begin{equation}\label{eq:rl-aufstieg}
  \boxed{\;P(X \text{ steigt auf}) = P_M(X) \cdot \sum_{Y} P_M(Y)\cdot P(X \text{ setzt sich gegen } Y \text{ durch})\;}
\end{equation}
Dabei ist $P_M$ die Meisterwahrscheinlichkeit aus dem Hauptlauf der
Staffel, demselben, aus dem auch die Spalte \enquote{Meister} der
Aufstiegsseite stammt, und $Y$ läuft über die Mannschaften der anderen
Staffel. Einen Malus-Lauf gibt es für die Regionalligen nicht
(Abschnitt~\ref{sec:malus}). Weil beide Größen aus einem Lauf kommen, gilt
immer $P(\text{Aufstieg}) \le P(\text{Meister})$.

\paragraph{Ein Duell über zwei Spiele.} Die Quote
$P(X \text{ setzt sich gegen } Y \text{ durch})$ rechnet
\datei{RCode/aufstiegsspiele.R:43--117} analytisch aus:

\begin{enumerate}
  \item Die Torerwartungen beider Spiele liefert der Endpunkt
        \code{/match-preview}, mit der \emph{aktuellen} ELO beider
        Mannschaften. Gewertete Spiele bewegen sie auch hier nicht
        (\code{rl\_aktuelle\_elo} sendet \code{elo\_neutral}, Issue~\#293).
        Im Hinspiel hat Nord Heimrecht, im Rückspiel Bayern.
  \item Die Tordifferenz eines Spiels ist verteilt als
        \[
          P(D = d) = \sum_{i} \pi_{\lambda_X}(i)\,\pi_{\lambda_Y}(i - d),
        \]
        abgeschnitten bei 15 Toren je Mannschaft.
  \item Die Gesamtdifferenz ist die Faltung $D_1 * D_2$ der beiden Spiele.
        Eine Auswärtstorregel gibt es nicht mehr.
  \item Bei Gleichstand folgt eine Verlängerung mit den \emph{neutralen}
        Raten $\lambda^{\text{neutral}} = (\lambda^{\text{hin}} + \lambda^{\text{rück}})/2$,
        auf 30 Minuten skaliert: $\lambda_{\text{ev}} = \lambda^{\text{neutral}}/3$.
        Ist es danach immer noch gleich, gilt das Elfmeterschießen als
        50:50.
\end{enumerate}
\begin{equation}
  P(X) = P(D_1 + D_2 > 0) + P(D_1 + D_2 = 0)\,\Bigl[P(D_{\text{ev}} > 0) + \tfrac12\,P(D_{\text{ev}} = 0)\Bigr].
\end{equation}
Im affinen Tormodell hebt sich der Heimvorteil über Hin- und Rückspiel
exakt auf. Wer zuerst Heimrecht hat, spielt für die Quote keine Rolle.

\paragraph{Siegquote.} Die Aufstiegsseite zeigt für die Playoff-Staffeln
zusätzlich die Siegquote
\[
  \frac{P(\text{Aufstieg})}{P(\text{Meister})},
\]
also die Chance, die Aufstiegsspiele zu gewinnen, wenn man Meister wird. Sie
ist die über die möglichen Gegner gemittelte Siegchance aus
Gleichung~\eqref{eq:rl-aufstieg} und deshalb nie größer als 100~Prozent. Ist
$P(\text{Meister}) = 0$, bleibt die Zelle leer.

\section{Aufstiegsberechtigung und Aufstiegssperre}\label{sec:malus}

Zweite Mannschaften von Profivereinen dürfen nicht in jede Liga aufsteigen.
Nach §\,55b Nr.~3.1 der DFB-Spielordnung gilt (ADR~0008, ADR~0009):

\begin{itemize}
  \item In der 3.~Liga und der 2.~Frauen-Bundesliga sind Zweitvertretungen
        grundsätzlich nicht aufstiegsberechtigt.
  \item In der Regionalliga ist eine Zweitvertretung nur dann nicht
        aufstiegsberechtigt, wenn die erste Mannschaft im kommenden Spieljahr
        in der 3.~Liga spielt (\emph{Aufstiegssperre}).
\end{itemize}

Dazu kommt der \emph{Aufstiegsverzicht}: Ein Verein kann für eine Saison auf
den Aufstieg verzichten (TeamList-Spalte \code{Aufstiegsverzicht}).

\paragraph{Masken statt zweitem Lauf.} Das Modell simuliert jede Liga
\emph{einmal}. Die Engine wertet dieselbe Abschlusstabelle zusätzlich unter
einer oder mehreren \emph{Berechtigungsmasken} aus und zählt, auf welchem Rang
\emph{unter den Berechtigten} jedes Team steht. Ein nicht berechtigter Meister
vererbt seinen Aufstieg an den nächsten Berechtigten. Wer nicht berechtigt ist,
steht in der TeamList (\code{Promotion}, \code{Aufstiegsverzicht}) und wird
nicht aus dem Kürzel erraten. Heatmap und Abstieg kommen unverändert aus dem
Tabellenlauf. Bis Oktober 2026 rechnete das Modell stattdessen einen zweiten
Lauf, in dem Gesperrte mit $-50$ Punkten starteten; die Maske ersetzt ihn
(ADR~0009).

\paragraph{Dynamische Sperre durch den Abstieg der Erstmannschaft.}
Spielt die Erstmannschaft einer Regionalliga-Reserve heute in der
2.~Bundesliga, steht erst am Saisonende fest, ob die Reserve gesperrt ist: Die
Erstmannschaft muss direkt absteigen oder die Relegation (Platz~16 gegen den
berechtigten Dritten der 3.~Liga) verlieren. Das Modell rechnet das exakt als
Mischung
\begin{equation}
  P(\text{Aufstieg}) = \sum_S P(S)\, M_S ,
\end{equation}
wobei $S$ die \emph{Menge} der gesperrten Reserven, $P(S)$ ihre gemeinsame
Verteilung und $M_S$ das Staffelergebnis unter der Maske \enquote{$S$ gesperrt}
ist. Das ist exakt, weil 2.~Bundesliga, 3.~Liga und die Staffeln kein Spiel
teilen. Die gemeinsame Verteilung zählt die Engine in der 2.~Bundesliga
(höchstens zwei Teams steigen direkt ab, die Reserven verschiedener Vereine
sind also gekoppelt), die Relegation des Sechzehnten löst R mit der
Gewinnwahrscheinlichkeit des Relegationsduells (zwei Spiele, Verlängerung mit
Heimvorteil des Rückspiel-Gastgebers) auf. Nord und
Bayern hängen an derselben 2.~Bundesliga; die Doppelsumme der Aufstiegsspiele
läuft deshalb über die gemeinsamen Gewichte beider Staffeln. Die Aufstiegsseite
nennt je betroffener Reserve die Wahrscheinlichkeit, dass sie gesperrt ist; die
Aufstiegszone an Platz~1 der Ligatabelle sinkt um die Wahrscheinlichkeit
eines gesperrten Meisters.

Welche Erstmannschaft zu welcher Reserve gehört, steht in den TeamList-Spalten
\code{Erstvertretung\_TeamID} und \code{Erstvertretung\_ShortText}; die API
liefert diese Zuordnung nicht.
